\phantomsection
\chapter{\lang{State of the Art and Related Work}{Estado da Arte e Trabalhos Relacionados}}
\phantomsection

\section{O Paradigma \textit{Local-First Software} e a Soberania dos Dados}

O modelo cliente-servidor é o paradigma histórico da internet há décadas. Em sua égide centralizada se consolidou o
modelo \textit{Cloud-First}, onde o estado principal da aplicação, os dados e outras informações dos usuários residem em servidores
centralizados, controlados por grandes provedores (empresas, bancos, governos, entre outros) que agem como a única fonte de
verdade daqueles dados (\textit{Single Source of Truth}). Embora esse paradigma tenha resolvido os problemas de colaboração
em tempo real e compartilhamento de dados, ele tornou as aplicações dependentes de conectividade e latência, além de colocar em risco a
privacidade dos dados do usuário.

Alternativamente à abordagem centralizada, \citeonline{kleppmann2019} propuseram um novo paradigma, guiado por sete ideais (velocidade, uso em
múltiplos dispositivos, funcionamento \textit{offline}, colaboração, longevidade, privacidade e controle do usuário). O \textit{Local-First software} define
um modelo não apenas com funcionalidades \textit{offline}, mas com uma verdadeira inversão de hierarquia no armazenamento de dados. A
cópia primária dos dados deve residir no computador local do usuário; o servidor, quando existe, é reduzido a um papel secundário
de arquivamento, roteamento e auxílio às conexões ponto a ponto.

Conciliar a ``independência da rede'' com a ``colaboração contínua'' é uma batalha desafiadora em sistemas distribuídos, uma
vez que, \textit{offline}, os estados e ramificações divergem. \citeonline{kleppmann2019} apontam os CRDTs como a base mais promissora para
conciliar esses dois ``ideais'', pois garantem a SEC ao mesclar (\textit{merge}) alterações locais de forma determinística quando
a rede é restabelecida, ainda que existam desafios em aberto.

O paradigma \textit{Local-First} traz consigo uma discussão antiga sobre soberania de dados. No modelo estabelecido, o usuário tem, supostamente,
controle de seus dados, mas precisa se comunicar por uma API para acessá-los, um direito temporário de acesso, em muitos casos. Caso ocorra um
bloqueio de conta, o usuário perde o acesso, mas aqueles dados não são necessariamente deletados pelo provedor e, em incontáveis casos, são usados
como produto para terceiros \cite{zuboff2019}. A célebre frase ``Dados são o novo petróleo'' por Clive Humby \cite{palmer2006} ilustra a visão
das empresas sobre esse aspecto.

Ao transferir a base de dados em sua completude para o armazenamento local, o sistema confere ao usuário a verdadeira custódia da informação. No
contexto \textit{peer-to-peer}, isso é elevado ao grau máximo de soberania: a centralização da infraestrutura é removida e a comunicação se torna direta
entre as máquinas que detêm posse local desses dados, unificando a autonomia do armazenamento local com a topologia descentralizada de rede. Implementações
em cima das redes P2P, voltadas à soberania de dados, incluem as \textit{blockchains} e a Web3.

\section{\textit{Blockchain} e Web3: Consenso Global versus Convergência sem Consenso}

Cada evolução das redes P2P até o presente foi marcada por diferentes aplicações novas surgindo em decorrência de um novo problema resolvido, como, por exemplo,
o bittorrent após a inovação no envio em chuncks de arquivos. Neste contexto, nas décadas mais recentes, duas implementações se destacam:
aplicações financeiras (bitcoin e ethereum) e aplicações em web 3.0.

Satoshi Nakamoto publicou em 2008 o artigo do Bitcoin, a meta era criar um sistema financeiro sem bancos capaz de realizar transações financeiras de modo
quase anônimo. O maior desafio técnico dessa rede era o "problema do gasto duplo": se um usuário detém uma moeda e tenta enviá-la para a pessoa A e para a pessoa
B ao mesmo tempo, a rede P2P precisa ter um mecanismo infalível para decidir qual transação vale e qual deve ser anulada. Para endereçar essa questão,
o Bitcoin e do Ethereum usam a rede P2P para trocar mensagens (gossip), mas o estado real da rede não muda. Eles agrupam as transações em blocos e usam o
poder computacional (Proof-of-Work) ou apostas financeiras (Proof-of-Stake) para garantir quem deve escrever no próximo bloco, atingindo o consenso global
síncrono. Toda a rede precisa concordar com uma ordem total e estrita dos eventos. Todos os computadores do mundo processam exatamente a mesma fila de
informações.

Para o caso da Web 3.0, o modelo anterior é inviável. Uma aplicação colaborativa não pode esperar segundos em uma fila de operações a cada movimento do mouse ou
edição de um documento. Aqui entra a teoria os CRDTs, usados para, de forma deterministica, convergir a rede para o mesmo estado, sem a necessidade de consenso.
Se o Usuário A move um marcador para a coordenada $X$, e o Usuário B move o mesmo marcador para a coordenada $Y$ exatamente no mesmo milissegundo, a rede não
precisa travar e fazer uma votação global para decidir qual edição prevalece. O caso colaborativo usa a rede P2P de forma otimista, o computador aceita a mudança localmente
na hora (Local-First, zero latência) e, em segundo plano, envia a edição para os outros nós da rede.

ESSA SEÇÃO PRECISA DE UMA MELHORA DE COESÃO, IDEIA E ESTRUTURA! WORK IN PROGRESS.

\section{Protocolos Sociais Descentralizados: AT Protocol e Nostr}

Outro caso de uso moderno de protocolos descentralizados são os protocolos sociais. Foram pautados em uma contraestratégia ao modelo das
redes sociais da web 2.0, onde as grandes empresas de tecnologia centralizadas controlam, vendem e espionam dados dos usuários a todo momento. De forma
contrária, a indústria e a comunidade \textit{open-source} começaram a buscar protocolos onde o usuário é dono do seu grafo social e dos seus dados.

Atualmente, essas redes sociais estão entre as maiores redes distribuídas em uso. A literatura e a documentação mostram como sincronizar milhões de pequenos
eventos diários sem a utilização engessada de uma blockchain. Porém, na prática, alguns ainda dependem de servidores grandes, como o Bluesky \cite{kleppmann2024bluesky}. O problema que eles resolvem é o dos ``jardins murados'' (\textit{walled gardens}) da
Web 2.0 (como Twitter e Facebook), onde a empresa central é dona do seu grafo social e dos seus dados \cite{bernerslee2010}. Se o usuário for banido ou a empresa falir, ele perde tudo.

\subsection{Nostr: Notes and Other Stuff Transmitted by Relays}

O Nostr é um protocolo minimalista que concentra o poder no cliente local do usuário em detrimento dos grandes servidores centrais, não é necessário um \textit{login},
criação de conta ou exposição de \textit{e-mail}/senha a terceiros. A identidade do usuário é um par de chaves criptográficas (pública e privada), semelhante ao modelo de identidade
das redes P2P. Porém, o Nostr não é P2P: os clientes se conectam a servidores (\textit{relays}) que não conversam entre si. Esses servidores se tornam ``burros'' agindo apenas como espalhador de mensagens criadas e assinadas criptograficamente nos nós
 ``inteligentes'' (computador, aplicativo do celular, televisão inteligente) locais dos usuários. Se algum desses servidores bloquear o usuário por quaisquer razões, seus
dados não estão perdidos, basta se conectar a outro servidor para espalhar suas mensagens. Como a chave criptográfica pertence ao usuário, seus seguidores o
encontrarão na nova rota, desde que saibam em quais \textit{relays} ele publica \cite{nip65}. Os dados e a identidade pertencem puramente ao nó local do usuário \cite{nip01}.

\subsection{AT Protocol: Authenticated Transfer Protocol (Bluesky)}

O Authenticated Transfer Protocol foi desenvolvido pela equipe do Bluesky (projeto iniciado pelo Twitter em 2019 e independente desde 2021) com participação de pesquisadores de sistemas distribuídos, como Martin Kleppmann \cite{kleppmann2024bluesky}.
Surge como o protocolo por trás da rede social Bluesky, focado na portabilidade de contas e na separação entre quem guarda os dados e quem mostra o \textit{feed} para o usuário. Alternativamente a um
banco de dados centralizado de um grande provedor, usuários têm o direito de hospedar o seu próprio PDS (\textit{Personal Data Server}) em um servidor próprio, que precisa ficar sempre \textit{online}, para guardar seu estado pessoal.

Para organizar o repositório de cada usuário, são empregadas \textit{Merkle Search Trees} \cite{auvolattaiani2019}. Elas servem de base para CRDTs de estado e permitem que os indexadores da
rede comparem o estado de dois servidores rapidamente e troquem apenas os deltas (as diferenças ou novas postagens), sem precisar copiar o banco de dados inteiro.
Como o estado dos dados está separado do restante, os usuários têm o direito a escolher qual algoritmo irá organizar o seu \textit{feed} personalizado, resultando em uma forma
de soberania algorítmica \cite{kleppmann2024bluesky}.

\section{libp2p: A Modular Networking Stack for Peer-to-Peer Applications}

\subsection{Contexto e motivação}

Historicamente, aplicações \textit{peer-to-peer} (como Napster ou BitTorrent) possuíam camadas de rede monolíticas, com protocolos de transporte e mecanismos de descoberta de pares
rigidamente acoplados à lógica de suas aplicações. Isso criava sistemas inflexíveis, incapazes de se comunicar entre nodos usando diferentes meios
de transporte ou ambientes (navegador, servidor, smartphone).

A motivação por trás da libp2p foi criar uma \textit{stack} de rede modular agnóstica. A internet é um ambiente fora de controle, repleto de \textit{firewalls}, roteadores NAT e dispositivos
operando em camadas e contextos distintos. A libp2p foi imaginada para integrar esses diferentes sistemas, independentemente de seus protocolos, redes, sistemas operacionais e
ambientes distintos \cite{libp2p}. Para isso, ela é composta por uma coleção de módulos que podem ser empilhados e usados em conjunto, conforme a necessidade da aplicação.

\subsection{Funcionamento}

Diferente do modelo cliente-servidor tradicional, onde os nós são identificados por endereços de IP, na libp2p cada nó é identificado por um PeerID, derivado da sua chave criptográfica pública \cite{libp2ppeerids}.
O endereço de IP é tratado apenas como uma rota temporária (multiaddr) para chegar àquela identidade \cite{multiaddr}.

A biblioteca suporta diversos meios de transporte, como TCP, QUIC (sobre UDP), WebRTC (para \textit{browsers}) e WebSockets \cite{libp2pspecs}. Uma vez que
dois nós tenham estabelecido uma conexão segura, é permitido que múltiplos subprotocolos paralelos rodem sobre essa mesma conexão (multiplexação), sem interferir uns nos outros \cite{libp2pconnections}. A ferramenta também
resolve o clássico problema de encontrar nós, através de módulos intercambiáveis \cite{libp2pspecs}.

Já para o problema de NAT, a libp2p oferece três módulos: o AutoNAT, que detecta se o nó está atrás de um NAT \cite{libp2pautonat}; o Circuit Relay, que usa um outro nó como intermediário
da conexão \cite{libp2prelay}; e o DCUtR, que tenta transformar essa conexão intermediada em uma conexão direta, através de \textit{hole punching} \cite{libp2pdcutr}.

\subsection{Implementações e Casos de Uso}

A forte adoção da libp2p validou o modelo modular para toda uma indústria. Ela atua como infraestrutura de transporte em grandes redes como a camada de consenso do Ethereum \cite{ethp2p} e o IPFS \cite{libp2p}.

\section{Iroh}

\subsection{Contexto e motivação}

Embora a libp2p e o IPFS (para sistemas de arquivos) tenham se estabelecido como padrões na indústria para redes \textit{peer-to-peer} e Web3, essas ferramentas ainda são consideradas, para algumas
aplicações, complexas. Como alternativa, surgiu uma segunda opção, que evoluiu conceitualmente e respondeu a gargalos da libp2p, com foco em velocidade, simplicidade e dispositivos móveis.
O acoplamento entre os componentes da libp2p (DHT, \textit{pubsub}, negociação de transporte) torna o sistema difícil de configurar e de entender \cite{irohfaq}.

O Iroh nasceu como um projeto de reescrita do IPFS na linguagem Rust para torná-lo mais rápido. Contudo, para atingir as métricas de desempenho e confiabilidade desejadas, os
desenvolvedores perceberam que precisariam quebrar a compatibilidade com as implementações existentes do IPFS \cite{irohnewdirection}. A motivação do Iroh é fornecer conexões P2P diretas entre dispositivos, identificados por chave pública, com a
menor latência possível e sem a complexidade de configuração da camada de rede \cite{iroh}.

\subsection{Funcionamento}

O Iroh foi construído nativamente sobre o protocolo QUIC \cite{rfc9000}, a base do HTTP/3 \cite{rfc9114}, descartando o TCP \cite{iroh}. O QUIC oferece estabelecimento rápido de conexão, criptografia nativa (TLS 1.3) e ausência do bloqueio de cabeça de fila (\textit{head-of-line blocking}) entre \textit{streams},
tornando a troca de mensagens mais rápida, mesmo em redes instáveis, como 4G e 5G.

Diferente de protocolos antigos que usam SHA-256, o iroh-blobs, protocolo de transferência de dados do Iroh, utiliza o algoritmo de hash BLAKE3 \cite{irohblobs}. Por ser internamente uma árvore de Merkle, ele permite que os dados sejam processados de forma paralelizada e verificados durante a transmissão (\textit{verified streaming}), aproveitando a velocidade
máxima do hardware \cite{blake3}.

Ele resolve problemas com NAT (roteadores bloqueando conexões de fora), um dos principais de aplicações P2P, com servidores de repasse (\textit{relays}), originalmente chamados de DERP (\textit{Designated Encrypted Relay for Packets}) \cite{irohchangelog}. A conexão começa passando pelo \textit{relay}, por onde os nós trocam informações de rede e tentam se conectar diretamente (\textit{hole punching}). Se
conseguirem, o \textit{relay} sai do caminho; se os \textit{firewalls} forem muito rígidos, o tráfego continua passando por ele. O \textit{relay} não consegue ler os dados, pois a conexão é criptografada de ponta a ponta \cite{irohrelay}. O usuário não
percebe a diferença, pois ela é abstraída pelo Iroh.

Sobre essa camada de conexão, o Iroh oferece protocolos separados, como o iroh-docs, para sincronização de estados \cite{irohprotocols}. Ele permite que múltiplos nós editem uma estrutura de dados de
chave-valor concorrentemente e a rede cuida de sincronizar as atualizações, através de uma técnica chamada \textit{range-based set reconciliation} \cite{irohdocs, meyer2022}.

\subsection{Implementações e Casos de Uso}

O Iroh já é utilizado em aplicações reais. O Delta Chat, um aplicativo de mensagens descentralizado \cite{deltachat}, usa o Iroh em duas funcionalidades: para transferir o
\textit{backup} de um dispositivo para outro através de um QR code \cite{deltachattransfer}, e para criar canais de comunicação em tempo real entre os aplicativos webxdc
que rodam dentro das conversas, através do iroh-gossip \cite{deltachatpeerchannels}. Nos dois casos, os dispositivos se conectam diretamente, sem que os dados passem pelo servidor de mensagens.

Já a Nous Research utiliza o Iroh na Psyche, uma rede para treinamento distribuído de modelos de IA pela internet. Os nós trocam metadados sobre os resultados do treinamento
através do iroh-gossip e baixam os resultados em si através do iroh-blobs. Segundo os autores, o Iroh consegue estabelecer conexões diretas em cerca de 90\% dos casos, contra
70\% da libp2p, e migra a conexão entre redes sem interrompê-la, graças ao QUIC \cite{nouspsyche}.

Por fim, os desenvolvedores do Iroh mantêm uma integração de exemplo com o Automerge, onde o Iroh é usado como camada de rede para sincronizar documentos CRDT entre os nós \cite{irohautomerge}.

\section{Yjs}

\subsection{Contexto e motivação}

A prova acadêmica dos CRDTs resolveram os problemas matemáticos para sincronização de estados, mas consumiam muita memória e CPU no processo. A serialização de grandes JSONs travava navegadores
mais fracos, a depender do tamanho dos documentos serializados. O Yjs foi criado por Kevin Jahns para preencher a lacuna entre a teoria acadêmica e alta performance, implemetando uma biblioteca capaz
de suportar a edição colaborativa em grande escala sem gargalos de memória e CPU.

Yjs implementa o algoritmo YATA, anteriormente desenhado para evitar o embaralhamento de textos quando usuários diferentes editam os mesmos valores. Ele modela o documento colaborativo como uma lista
duplamente encadeada, em vez de árvores. O Run-Length Encoding (RLE) usa codifiação binária compacta para condensar milhares de operações em vetor de poucos bytes. O Yjs é focado nos dados e agnóstico
em relação à rede, permitindo acoplar os deltas binários em WebRTC, WebSockets ou malhas P2P.

\subsection{Funcionamento}

O Yjs não faz suposições sobre a rede usada nas implementações. Paradigmas como cliente-servidor são tão usáveis quanto \textit{peer-to-peer}: desde que as mudanças sejam comunicadas, ocorrerá a sincronização dos estados.
Isso ocorre independentemente de a comunicação ser feita por WebSockets, WebRTC (\textit{browsers}) ou outros \textit{providers} \cite{yjs}. Ele suporta edição \textit{offline}, o que permite seguir a filosofia Local-First, onde as alterações são aplicadas
primeiramente no local e depois comunicadas aos outros nodos. Eles, por sua vez, aplicam essas alterações, que são mescladas automaticamente pelos CRDTs, abstraídos como estruturas de dados convencionais do JavaScript.

Internamente, todos os tipos compartilhados (\textit{shared types}) são listas. Isso significa que um único algoritmo, o YATA \cite{nicolaescu2016}, resolve os conflitos de todos os tipos. Cada item tem um ID, formado pelo identificador do cliente e por um relógio de Lamport, e guarda os vizinhos
do momento em que foi inserido. Com isso, dois usuários que inserem no mesmo ponto chegam à mesma ordem final. Um item deletado não desaparece, ele vira um \textit{tombstone} no conjunto de deleções (\textit{delete set}). Para que os metadados não cresçam \textit{ad infinitum},
o Yjs junta em um único item os caracteres digitados em sequência pelo mesmo usuário e descarta o conteúdo dos itens deletados. Em um documento real, com 182 mil inserções e 77 mil deleções, o conjunto de deleções ocupou apenas 4,5 KB \cite{yjsinternals}.

\subsection{Implementações e Casos de Uso}

Usado por diversas aplicações \cite{yjs}, o Yjs é utilizado tanto em aplicações cliente-servidor quanto \textit{peer-to-peer}. Entre os usuários mais conhecidos estão o JupyterLab, para notebooks colaborativos, e o JupyterGIS, um editor colaborativo de sistemas de informação geográfica \cite{yjs, jupytergis}. O Room.sh, uma ferramenta de reuniões com funcionalidades embutidas de
desenho, edição e codificação, faz uso do \textit{framework} para sincronizar seus vários estados entre múltiplos usuários. Outras aplicações diferentes também contam com uso do Yjs, como o
Dynaboard, um construtor colaborativo de aplicações web, e o LegendKeeper, uma ferramenta colaborativa para planejamento de campanhas de RPG \cite{yjs}.

\section{Automerge}

\subsection{Contexto e motivação}

A motivação do Automerge era permitir aos programadores a criação de softwares Local-First manipulando estruturas de dados familiares (listas, \textit{strings}, objetos JSON) sem precisar lidar com os problemas complexos de sistemas distribuídos \cite{automerge}. Nascido com foco na experiência do desenvolvedor e na preservação de histórico, o Automerge
funciona como um ``Git para estruturas de dados'' \cite{automerge}. Foi criado pelo laboratório Ink \& Switch, o mesmo que publicou o artigo fundador do Local-First, com participação de Martin Kleppmann, um dos autores desse artigo, e é baseado na pesquisa sobre CRDTs para JSON \cite{kleppmann2019, kleppmannberesford2017}.

\subsection{Funcionamento}

Enquanto o Yjs representa internamente todos os tipos como listas, com foco em desempenho, o Automerge modela o documento como uma árvore JSON aninhada \cite{kleppmannberesford2017}. Na versão 3.0, o Automerge passou a usar o seu formato colunar comprimido também durante a execução, e não apenas no armazenamento, reduzindo o uso de memória em mais de 10 vezes \cite{automerge3}.
Ele salva o histórico completo das mutações compactando colunas de dados semelhantes \cite{automerge3}. Já a sincronização envia apenas as alterações que o outro nó ainda não possui \cite{automerge}.

Assim como no Git, cada alteração feita em um documento é registrada como uma \textit{change}, parecida com um \textit{commit}. Cada \textit{change} é identificada pelo
\textit{hash} SHA-256 do seu conteúdo e aponta para as \textit{changes} das quais depende, formando um grafo acíclico dirigido (DAG). As \textit{changes} mais recentes desse grafo,
que ainda não possuem sucessoras, são chamadas de \textit{heads} \cite{automergespec}. Dentro de cada \textit{change}, toda operação recebe um ID formado por um contador e
pelo identificador da réplica (\textit{actor ID}), ou seja, um relógio de Lamport. Isso impõe uma ordem total entre as operações: vence o maior contador e, em caso de empate,
o maior \textit{actor ID} \cite{automergespec}.

Na maioria dos casos, a mesclagem é direta. Alterações em objetos ou propriedades diferentes são simplesmente combinadas. Em listas e textos, todas as inserções e deleções
são preservadas, e inserções concorrentes na mesma posição recebem a mesma ordem em todos os nós \cite{automergeconflicts}. Já para valores numéricos, o Automerge oferece
o tipo \texttt{Counter}, que soma incrementos concorrentes: se dois usuários incrementam um valor 3 ao mesmo tempo, o resultado é 5, e não 4 \cite{automergecounters}.

O único caso sem resolução automática é quando dois usuários alteram a mesma propriedade do mesmo objeto. Nesse caso, o Automerge escolhe um dos valores como vencedor,
de forma determinística e igual em todos os nós, sem depender do relógio físico das máquinas. Os valores perdedores não são descartados: eles ficam acessíveis pela função
\texttt{getConflicts()}, permitindo que a aplicação mostre o conflito ao usuário \cite{automergeconflicts}.

A sincronização entre dois nós é baseada no trabalho de \citeonline{kleppmannhoward2020}. Cada nó guarda o que sabe sobre as \textit{heads} do outro e, a cada mensagem
trocada, envia apenas as \textit{changes} que o outro ainda não possui, até que os dois tenham o mesmo histórico \cite{automergesync}. Por cima disso, o Automerge Repo
cuida da parte de infraestrutura: ele salva os documentos localmente através de adaptadores de armazenamento, como o IndexedDB no navegador ou o sistema de arquivos,
e os distribui para outros nós através de zero ou mais adaptadores de rede, como o WebSocket \cite{automergerepo}.

\subsection{Implementações e Casos de Uso}

Os casos de uso podem variar, mas ele se estende para todas as aplicações Local-First ou que requerem um histórico de edições, como no Git. Aplicações Local-First exigem armazenamento local robusto e controle de versão, como gerenciadores de conhecimento, ferramentas de produtividade \textit{offline} e outros tipos de sistemas
onde os usuários necessitam viajar entre os históricos passados de alterações.

Os primeiros casos de uso foram os protótipos do próprio Ink \& Switch, apresentados no artigo sobre Local-First \cite{kleppmann2019}: o Trellis, um quadro Kanban inspirado no
Trello, que experimentou uma interface de ``viagem no tempo'' para navegar pelos estados anteriores do documento; o Pixelpusher, um editor colaborativo de \textit{pixel art}; e o
PushPin, um mural colaborativo de mídias mistas, parecido com o Miro, considerado pelos autores o mais completo dos três. Além disso, como visto na seção sobre o Iroh, existe
uma integração de exemplo que usa o Iroh como camada de rede para sincronizar documentos do Automerge \cite{irohautomerge}.

\section{OrbitDB e Merkle-CRDTs}

\subsection{Contexto e motivação}

A libp2p é usada como a camada de rede do protocolo IPFS \cite{libp2p}. O IPFS é excelente para distribuir arquivos estáticos de forma imutável, mas não funciona bem para arquivos que mudam dinamicamente, como em um banco de dados,
pois os arquivos são endereçados pelo \textit{hash} do seu conteúdo e qualquer alteração gera um novo identificador \cite{benet2014}.
O OrbitDB foi criado para ser um banco de dados distribuído e sem servidor, voltado para aplicações P2P, Web3 e Local-First. Ele une o armazenamento P2P do IPFS com a consistência eventual dos CRDTs, eliminando completamente a figura do servidor de banco de dados tradicional \cite{orbitdb}.

\subsection{Funcionamento}

Todo banco de dados no OrbitDB é formado por um registro de operações (\textit{oplog}), um \textit{log} imutável onde novas operações são apenas adicionadas ao final (\textit{append-only}). O estado do banco é
calculado a partir dessas operações. Sobre o \textit{oplog}, o OrbitDB oferece diferentes tipos de banco de dados: \textit{events} (um \textit{log} de eventos), \textit{documents} (documentos JSON indexados por uma chave),
\textit{keyvalue} e \textit{keyvalue-indexed} (chave-valor) \cite{orbitdb}.

O \textit{oplog} é formalmente um Merkle-CRDT \cite{sanjuan2020}. Cada operação no banco de dados, como adicionar ou remover uma chave, vira uma entrada imutável armazenada no IPFS, identificada pelo seu \textit{hash} criptográfico.
Cada entrada referencia as entradas anteriores na ordem causal, e o banco em sua completude forma um Grafo Acíclico Dirigido (DAG), semelhante aos \textit{commits} do Git. Como cada entrada é identificada pelo \textit{hash}
do seu conteúdo e das suas referências, nenhuma entrada pode ser alterada ou removida sem que isso seja detectado \cite{orbitdboplog}. A ideia central dos Merkle-CRDTs é que o próprio DAG funciona como um relógio lógico, o que
simplifica a construção de CRDTs em redes onde as mensagens podem ser perdidas, duplicadas ou chegar fora de ordem \cite{sanjuan2020}.

A sincronização é feita através da libp2p. Ao abrir um banco, o nó se inscreve em um tópico de \textit{pubsub} com o identificador do \textit{log}. Os nós que já estão no tópico se conectam diretamente ao novo nó e os dois trocam
as \textit{heads} do \textit{log}. Depois disso, a cada nova operação, o nó publica as suas novas \textit{heads} no tópico, e os outros nós buscam as entradas que ainda não possuem. O protocolo não garante a ordem nem a entrega das mensagens,
mas garante que, quando todas forem entregues, todos os nós terão o mesmo estado, ou seja, consistência eventual \cite{orbitdbsync}.

\subsection{Implementações e Casos de Uso}

A implementação principal do OrbitDB é em JavaScript e funciona tanto no navegador quanto no Node.js. Existe também uma implementação em Go, a go-orbit-db, desenvolvida e mantida pelo projeto Berty \cite{orbitdb}. O Berty é um
aplicativo de mensagens \textit{offline-first}, \textit{peer-to-peer} e sem servidor \cite{berty}. No seu protocolo, chamado Wesh, as conversas acontecem em ``grupos'', onde múltiplos dispositivos compartilham mensagens e metadados através do OrbitDB \cite{berty}.

Os Merkle-CRDTs também são usados fora do OrbitDB. O go-ds-crdt, um banco chave-valor baseado no mesmo artigo de \citeonline{sanjuan2020}, replica automaticamente cada par chave-valor para todos os nós da rede, mesmo que as mensagens
sejam perdidas, reordenadas ou duplicadas \cite{godscrdt}. Ele é usado pelo IPFS Cluster, uma ferramenta que coordena quais arquivos devem ser replicados entre vários nós do IPFS, para manter esse estado global sincronizado
entre todos os nós do \textit{cluster} \cite{ipfscluster}.

\section{Sincronização Colaborativa de Dados Espaciais (WebGIS)}

WORK IN PROGRESS - É PRECISO AVANÇAR MAIS PARA CONTINUAR ESSE CAPITULO

\section{Comparativo e Posicionamento do Trabalho}

WORK IN PROGRESS - É PRECISO AVANÇAR MAIS PARA CONTINUAR ESSE CAPITULO

\begin{table}[htb]
    \centering
    \caption{Comparativo entre libp2p e Iroh}
    \label{tab:libp2p-iroh}
    \begin{tabular}{|p{0.22\textwidth}|p{0.34\textwidth}|p{0.34\textwidth}|}
        \hline
        \textbf{Critério} & \textbf{libp2p} & \textbf{Iroh} \\ \hline
        Arquitetura & Coleção de módulos empilháveis, conforme a necessidade da aplicação & Núcleo de conexão, com protocolos separados por cima (iroh-blobs, iroh-gossip, iroh-docs) \\ \hline
        Implementações & Várias (Go, Rust, JavaScript, entre outras), seguindo as mesmas especificações & Rust \\ \hline
        Identidade do nó & PeerID, derivado da chave pública & Chave pública Ed25519 \\ \hline
        Transportes & TCP, QUIC, WebRTC, WebSockets, entre outros & QUIC \\ \hline
        Segurança e multiplexação & Módulos próprios sobre a conexão (ou nativos, no QUIC) & Nativos do QUIC \\ \hline
        Travessia de NAT & AutoNAT, Circuit Relay e DCUtR (\textit{hole punching}) & Conexão inicia pelo \textit{relay} e migra para direta via \textit{hole punching} \\ \hline
        Conexões diretas (segundo a Nous Research) & Cerca de 70\% & Cerca de 90\% \\ \hline
        Descoberta de pares & Kademlia DHT, mDNS e \textit{rendezvous} & DNS (Pkarr) e, opcionalmente, a DHT do BitTorrent; sem DHT Kademlia própria \\ \hline
        Difusão de mensagens & GossipSub & iroh-gossip (HyParView e PlumTree) \\ \hline
        Casos de uso & Camada de consenso do Ethereum, IPFS, OrbitDB & Delta Chat, Psyche \\ \hline
    \end{tabular}
    \fonte{Elaborado pelo autor, com base em \citeonline{libp2p}, \citeonline{libp2pspecs}, \citeonline{iroh}, \citeonline{irohfaq}, \citeonline{irohaddresslookup} e \citeonline{nouspsyche}.}
\end{table}

\begin{table}[htb]
    \centering
    \caption{Comparativo entre Yjs e Automerge}
    \label{tab:yjs-automerge}
    \begin{tabular}{|p{0.26\textwidth}|p{0.32\textwidth}|p{0.32\textwidth}|}
        \hline
        \textbf{Critério} & \textbf{Yjs} & \textbf{Automerge} \\ \hline
        Base teórica & YATA \cite{nicolaescu2016} & CRDT para JSON \cite{kleppmannberesford2017} \\ \hline
        Modelo de dados & Todos os tipos são listas internamente & Documento JSON aninhado (mapas, listas, texto, contadores) \\ \hline
        Histórico & Descarta o conteúdo dos itens deletados & Guarda o histórico completo, como um grafo de \textit{changes} \\ \hline
        Foco & Edição colaborativa em tempo real e desempenho & Versionamento e aplicações Local-First \\ \hline
        Rede & Agnóstico, via \textit{providers} & Agnóstico, via Automerge Repo \\ \hline
        Tempo para aplicar o histórico real (B4) & 5.714 ms & 14.326 ms \\ \hline
        Tamanho final do documento (B4) & 159.929 bytes & 129.116 bytes \\ \hline
    \end{tabular}
    \fonte{Elaborado pelo autor, com base em \citeonline{yjs}, \citeonline{yjsinternals}, \citeonline{automerge} e \citeonline{crdtbenchmarks}.}
\end{table}
